\documentclass[11pt]{article}
\usepackage[margin=1in]{geometry}
\usepackage{amsmath,amssymb}
\usepackage{enumitem}
\newcommand{\warn}{\textbf{[!]}}

\title{ME34: Mirrors \& Lenses --- Walkthrough}
\author{}\date{}

\begin{document}\maketitle

\warn{} \emph{No source answer key. Q1 inline answers reproduced from .tex; Q8 depends on figure origin.}

\section*{Sign convention}
\begin{itemize}[leftmargin=*]
\item $\dfrac{1}{s}+\dfrac{1}{s'} = \dfrac{1}{f}$ (mirror or thin lens)
\item $s>0$ real object, $s'>0$ real image, $f>0$ converging lens / concave mirror.
\item $m = -s'/s$
\item Refraction at spherical surface: $\dfrac{n_1}{s}+\dfrac{n_2}{s'} = \dfrac{n_2-n_1}{R}$
\item Lensmaker: $\dfrac{1}{f} = (n-1)\!\left(\dfrac{1}{R_1}-\dfrac{1}{R_2}\right)$
\end{itemize}

\section*{Q1 --- Flat mirror, $s=21$}
\textbf{Setup.} A plane (flat) mirror is the simplest case: an object 21 cm in front always forms an image the same distance behind the mirror, same size, upright, and virtual --- no real light converges there, your brain just traces reflected rays back.

(a) $-21$ cm; (b) $m=+1$; (c) virtual; (d) upright.

\section*{Q2 --- Concave mirror, $f=10.7$, $s=23$}
\textbf{Setup.} A concave (converging) mirror with $f=10.7$ cm and the object at $s=23$ cm $> f$: object is outside the focal point, so we expect a real, inverted image on the same side. Plug into $1/s+1/s'=1/f$.

$1/s' = 1/10.7 - 1/23 = 0.04998 \Rightarrow s' = \boxed{20.0\ \text{cm}}$ (real, inverted).

\section*{Q3 --- Convex mirror, $R=5.8$, $s=1.3$}
\textbf{Setup.} A convex (diverging) mirror has $f = -R/2 = -2.9$ cm. Convex mirrors always form virtual, upright, reduced images, so we expect $0<m<1$.

Convex: $f = -2.9$. $1/s' = -1/2.9 - 1/1.3 = -1.114 \Rightarrow s' = -0.898$, $m = -s'/s = \boxed{+0.690}$.

\section*{Q4 --- Converging lens, $f=14$, $s=10$ (object inside $f$)}
\textbf{Setup.} Converging lens but object \emph{inside} the focal length. Rays don't converge on the far side; traced back they appear to diverge from a magnified, upright virtual image on the same side as the object --- the magnifying-glass regime. Expect $s'<0$, $m>0$, $|m|>1$.

$1/s' = 1/14 - 1/10 = -0.02857 \Rightarrow s' = \mathbf{-35\ \text{cm}}$. $m = +3.5$. \textbf{Virtual, upright.}

\section*{Q5 --- Diverging lens, $f=-9$, $s=16$}
\textbf{Setup.} A diverging lens always produces a virtual, upright, reduced image of a real object. Expect $s'<0$ and $0<m<1$.

$f=-9$, $s=16$. $1/s' = -0.1736 \Rightarrow s' = \mathbf{-5.76\ \text{cm}}$, $m = \mathbf{+0.36}$. \textbf{Virtual, upright.}

\section*{Q6 --- Lensmaker, then image position}
\textbf{Setup.} A high-index ($n=2$) thin lens with $R_1=8$ cm and $R_2=-11$ cm (both surfaces convex outward, sign convention reflected in $R_2<0$). First compute $f$ from the lensmaker equation, then locate the image of an object at $s=23$ cm via the thin-lens equation.

$1/f = (2-1)(1/8 - 1/(-11)) = 0.2159 \Rightarrow f = 4.63$ cm.
$1/s' = 0.2159 - 1/23 = 0.1724 \Rightarrow s' = \boxed{5.80\ \text{cm}}$.

\section*{Q7 --- Pebble in sphere}
\textbf{Setup.} A pebble is embedded in a transparent sphere ($n=2$, $R=33$ cm); its light refracts once at the curved sphere-air boundary. The image you see appears inside the sphere (virtual). Use the single-surface refraction equation with careful signs: $s>0$ (real object inside), $s'<0$ (virtual image, same side), and $R<0$ because the center of curvature is on the incoming side from the standpoint of light heading outward.

$n_1=2$ (inside), $n_2=1$ (outside). Image inside (virtual): $s' = -14$. Center of curvature inside: $R = -33$.
\[ \frac{2}{s} + \frac{1}{-14} = \frac{1-2}{-33} \Rightarrow \frac{2}{s} = \frac{1}{33}+\frac{1}{14} = \frac{47}{462} \Rightarrow s = \boxed{19.66\ \text{cm}} \]

\section*{Q8 --- Two-lens system}
\textbf{Setup.} Tree at 21 cm in front of converging lens 1 ($f_1=+8$), with diverging lens 2 ($f_2=-17$) at $L=10$ cm down the axis. Solve sequentially: image from lens 1 becomes the object for lens 2. If lens 1's image would form past lens 2, it's a \emph{virtual} object for lens 2 with $s_2<0$.

Lens 1: $1/s_1' = 1/8 - 1/21 = 0.07738 \Rightarrow s_1' = +12.92$ cm.
Lens 2 (separation 10 cm): $s_2 = -(12.92-10) = -2.92$ (virtual object).
$1/s_2' = 1/(-17) - 1/(-2.92) = 0.2837 \Rightarrow s_2' = +3.53$ cm right of lens 2.
With origin at the object: $x_{\text{img}} = 21 + 10 + 3.53 = \boxed{34.53\ \text{cm}}$.

\section*{Q9 --- Corrective lens for myopia}
\textbf{Setup.} The patient's far point is 43.3 cm; everything past that is blurry. The corrective lens takes light from infinity and forms a virtual image at her far point so her eye can focus it. Use $s\to\infty$, $s'=-0.433$ m; lens power $P=1/f$ in diopters with $f$ in meters. Diverging lens for myopia, so $P<0$.

$f = 1/P = -0.4333$ m: $P = 1/(-0.433) = \boxed{-2.31\ \text{D}}$

\section*{Q10 --- Magnifier focal length}
\textbf{Setup.} Simple magnifier with target angular magnification $M=2.3$. Standard relaxed-eye formula (image at infinity) is $M = 25\ \text{cm}/f$; alternate near-point convention gives $M=1+25/f$. Use whichever your textbook does.

$f = 25/M = 25/2.3 = \boxed{10.87\ \text{cm}}$ (relaxed eye); or $f = 25/(M-1) = 19.23$ cm (image at near point).

\section*{Q11 --- Corrective optics}
Near-sightedness: eye focuses in front of the retina, so a \textbf{diverging} lens spreads incoming rays first; it forms a virtual image of distant objects at the eye's far point. Far-sightedness: eye can't focus close enough, so a \textbf{converging} lens forms a virtual image of nearby print farther away (at the near point) where the eye can focus.

Near-sightedness: \textbf{diverging} lens, image \textbf{virtual}. Far-sightedness: \textbf{converging} lens, image \textbf{virtual}.

\section*{Q12 --- Spherical mirror image type}
For a concave mirror, the focal point splits the two regimes. Object beyond $f$ gives a real, inverted image (rays converge on the same side as the object). Object inside $f$ gives a virtual, upright, magnified image (rays diverge after reflection; traced back, they form the image behind the mirror).

Object farther than $f$ (concave mirror): \textbf{real, inverted}. Closer than $f$: \textbf{virtual, upright}.

\bigskip\warn{} Q8's $x$-position depends on figure origin.
\end{document}
